\section{Product norm controls Boolean description size independently of depth}
\label{sec:boolean-product-depth-saturation}

The rank-one bottleneck strategy of Golowich--Rakhlin--Shamir
\cite{golowich2018} underlies depth-independent norm-based complexity
bounds. Here we use it for a more specific representation statement:
global Boolean margin permits a short, polynomial-norm tanh
implementation of the scalar suffix, without carrying a constant rail
through the retained prefix. The conclusion bounds a finite Boolean
description independently of the original depth as well as width.

\begin{lemma}[Relevant coordinates and freezing]
\label{lem:pf-junta-freezing}
A Boolean function with a globally unit-margin tanh representation
of product norm at most $P$ depends on at most $P^2$ coordinates.
After deleting all irrelevant coordinates, it has a representation
of the same depth and no larger product norm.
\end{lemma}

\begin{proof}
Let $T=\prod_{j\ge2}\|W^j\|_F$. For every relevant coordinate $i$,
some cube edge in that coordinate has opposite labels at its endpoints.
The two outputs differ by at least two, whereas the Lipschitz bound
along that edge is $2\|W^1e_i\|_2T$. Thus
$\|W^1e_i\|_2T\ge1$. Squaring and summing over the relevant coordinates
gives their number at most $\|W^1\|_F^2T^2\le P^2$.

Freeze each irrelevant coordinate $i$ to an independent uniform sign
$a_i$. This replaces the first-layer bias column $w_0$ by
$w_0+\sum_i a_iw_i$, while retaining the relevant columns. Its expected
squared Frobenius norm is exactly the original squared Frobenius norm.
Some assignment therefore does not increase that norm. Since the
function is independent of these coordinates, the reduced network
retains its unit margin on the full reduced cube.
\end{proof}

\begin{lemma}[A bias-free odd scalar classifier]
\label{lem:odd-scalar-classifier}
Let $P,R\ge1$ and let $g:[-R,R]\to\mathbb R$ be odd and $P$-Lipschitz.
Let $T\subset[-R,R]$ be finite and nonempty, with
$|g(t)|\ge7/8$ for $t\in T$. There is a tanh network with two hidden
layers, a linear readout and no biases, whose sign agrees with $g$
on $T$ and whose signed output there is at least $7/4$.
Its width, weights, and product of Frobenius norms are bounded by
universal polynomials in $P,R$.
\end{lemma}

\begin{proof}
Extend the set to $T\cup(-T)$ and sort its positive absolute values.
All are at least $7/(8P)$, since $g(0)=0$. Consecutive values with
opposite signs of $g$ are separated by at least $7/(4P)$.
Place a threshold $\theta_j$ at the midpoint of each such transition
gap. There are $m\le2PR$ transitions, every $0<\theta_j\le R$,
and every actual absolute value is at distance at least $7/(8P)$
from every threshold.

The resulting odd step classifier has the exact expression
\[
 F(t)=c_0\operatorname{sign}(t)+
   \sum_{j=1}^m c_j\operatorname{sign}
       \bigl(t-\theta_j\operatorname{sign}(t)\bigr),
 \qquad |c_0|\le1,\quad c_j\in\{-1,1\}.
\]
Indeed, the positive-side label changes by $2c_j$ at transition $j$.
The changes alternate, so the coefficient
$c_0=\ell_0+\sum_jc_j$ is in $\{-1,0,1\}$, where $\ell_0$ is the
first positive-side label. The displayed formula is the odd extension
of this step classifier and equals $\operatorname{sign}g(t)$ on $T$.

Set
\[
 \rho=\frac1{2P},\qquad
 \delta=\min\left\{\frac\rho{4R},\frac1{32(m+1)}\right\},\qquad
 A=\frac1{2\rho}\log\frac2\delta,
\]
\[
 \alpha^2=\frac\rho{4R^3},\qquad
 B=\rho^{-1}\log(64(m+1)),\qquad
 c=\operatorname{arctanh}(1/2).
\]
The first hidden layer computes
$u(t)=\tanh(\alpha t)$ and $s(t)=\tanh(At)$.
On $T\cup(-T)$,
\[
 |u(t)/\alpha-t|\le\rho/12,
 \qquad |s(t)-\operatorname{sign}(t)|\le\delta.
\]
Thus replacing $t-\theta_j\operatorname{sign}(t)$ by
$u(t)/\alpha-\theta_js(t)$ changes it by at most
$\rho/12+R\delta\le\rho/3$. The original magnitude exceeds $\rho$,
so its sign is unchanged and its new magnitude exceeds $\rho/2$.
The second hidden layer computes
\[
 r(t)=\tanh(cs(t)),\qquad
 v_j(t)=\tanh\bigl(B[u(t)/\alpha-\theta_js(t)]\bigr).
\]
Consequently $2r(t)$ approximates $\operatorname{sign}(t)$ to error
at most $2\delta$, and each $v_j$ approximates its desired sign to
error at most $1/[32(m+1)]$. Therefore
\[
 z(t)=2\left[2c_0r(t)+\sum_{j=1}^m c_jv_j(t)\right]
\]
has signed value at least
$2[1-2\delta-m/(32(m+1))]\ge15/8>7/4$ on $T$.
This includes $m=0$, when only the sign rail is needed.

The three squared matrix norms are bounded respectively by
\[
 \alpha^2+A^2,\qquad
 c^2+mB^2(\alpha^{-2}+R^2),\qquad 16+4m.
\]
Since $m\le2PR$, these expressions, their product, and every physical
weight are polynomially bounded in $P,R$. All biases are zero.
\end{proof}

\begin{theorem}[A depth-independent Boolean normal form]
\label{thm:depth-independent-pf-normal-form}
There is a universal polynomial $p$ with the following property.
Suppose $f:\{-1,1\}^d\to\{-1,1\}$ has some finite-depth,
finite-width tanh representation with unit global margin and product
of Frobenius norms at most $P\ge1$. Then $f$ has a globally
unit-margin representation using only its relevant input coordinates,
with depth, hidden widths, individual weights, and product of
Frobenius norms at most $p(P)$. Its weights can be chosen dyadic
with at most $p(P)$ bits each. The bounds are independent of the
original depth, widths, individual weights and their precision.
There are at most $P^2$ relevant coordinates.
\end{theorem}

\begin{proof}
First use Lemma~\ref{lem:pf-junta-freezing}. If no coordinate is
relevant, a constant network of width one and depth two suffices.
Otherwise write $k\le P^2$ for the number of remaining coordinates
and $R=\sqrt{k+1}$. Let $D$ be the remaining original depth, and put
\[
 r=\left\lceil128P^2R^2\log(2PR)\right\rceil.
\]
If $D\le r+2$, retain the network for the moment. Otherwise we
produce a representation with at most $r+2$ weighted layers and
polynomial product norm.

Write $M_j=\|W^j\|_F$, $s_j=\|W^j\|_{\rm op}$ and
$\operatorname{sr}_j=M_j^2/s_j^2$. Unit margin implies
$\prod_js_j\ge1/R$, because all activations are zero-preserving
and one-Lipschitz. Hence
\[
 \sum_j\log\operatorname{sr}_j
 =2\log\frac{\prod_jM_j}{\prod_js_j}
 \le2\log(PR).
\]
Among the first $r$ hidden layers choose one, $j$, with
$\log\operatorname{sr}_j\le2\log(PR)/r$.
Replace $W^j$ by its leading singular component $s_juv^T$.
The discarded operator norm is at most
$M_j\sqrt{1-1/\operatorname{sr}_j}
 \le M_j\sqrt{\log\operatorname{sr}_j}$.
Propagation through the original prefix and suffix therefore changes
every output by at most
\[
 PR\sqrt{2\log(PR)/r}\le1/8.
\]
The modified network retains the target's signs and margin at least
$7/8$.

Put $S=\prod_{i<j}M_i$, with $S=1$ if $j=1$, and normalize the
bottleneck scalar as $t(x)=v^Th^{j-1}(x)/S$. Its absolute value is at
most $R$. Everything from the rank-one matrix onward is an odd
scalar function $g(t)$ with Lipschitz constant at most
$s_jS\prod_{i>j}M_i\le P$. Apply
Lemma~\ref{lem:odd-scalar-classifier} on the finite set of these
actual scalar values.

Attach its two-hidden-layer classifier to the original prefix.
The first two new rows are $\alpha v^T/S$ and $Av^T/S$.
Their common factor $1/S$ cancels the prefix product $S$ in the
total Frobenius product. Thus the complete new network has polynomial
product norm in $P,R$, independent of every original dimension and
weight scale. Its depth is $j+2\le r+2$, and it has the required
global signs with margin at least $7/4$.

In either case we now have depth and product norm polynomial in $P$,
while widths may still be arbitrary. Apply the finite-description
normal form, Corollary~\ref{cor:product-finite-description}, with
input dimension $k\le P^2$. It gives polynomial width, physical
weights and dyadic bit length, keeping unit margin and a polynomial
product norm. Restore zero columns for the irrelevant original
coordinates if a $d$-input network is desired. If the normal form is
purely affine, apply tanh to that affine output and append the
readout $2$; its unit margin is preserved and all bounds remain
polynomial, now within the usual architecture with a hidden layer.
\end{proof}

\begin{corollary}[Counting globally margined product-norm classifiers]
\label{cor:all-depth-pf-counting}
Let $\mathcal H_P$ contain all Boolean functions admitting any
finite-depth, finite-width globally unit-margin tanh representation
with product norm at most $P\ge1$. For a universal polynomial $p$,
\[
 \log|\mathcal H_P|\le p(P)+P^2\log(d+1).
\]
Consequently, with probability at least $1-\delta$, every member
of $\mathcal H_P$ consistent with $n$ independent realizable examples
has error at most $\varepsilon$, provided
\[
 n\ge\frac{p(P)+P^2\log(d+1)+\log(1/\delta)}{\varepsilon}.
\]
There is no depth or width restriction in this simultaneous event.
\end{corollary}

\begin{proof}
There are at most $(d+1)^{\lfloor P^2\rfloor}$ possible ordered
lists of relevant coordinates, allowing a padding symbol. Once the
coordinates are chosen, Theorem~\ref{thm:depth-independent-pf-normal-form}
encodes the entire polynomial-size dyadic network in $p(P)$ bits,
after enlarging a universal polynomial. Count these descriptions and
apply the usual union bound: an error-$\varepsilon$ classifier is
consistent with $n$ independent examples with probability at most
$e^{-n\varepsilon}$.
\end{proof}

In particular, counting implies that most Boolean functions of $k$
variables require product norm at least $2^{c k}$ for some universal
$c>0$, even when arbitrary depth and width are allowed. This is not
a small-circuit lower bound: those generic functions also require
exponentially large ordinary Boolean descriptions.
